\documentclass[a4paper,12pt,fleqn]{article}%fleqn pour aligner les equations à gauche
\usepackage{../../Styles/Cours}


\newcommand{\chnum}{Chapitre 1}
\newcommand{\chtitle}{Électrisation}
\newcommand{\dtype}{Activité 1}

\begin{document}
	\maketitle
 Dès l'Antiquité, on observe que l'ambre, après avoir été frotté contre de la laine acquiert la capacité d'attirer certains corps légers. Le terme "électricité" a été créé au {\romannumeral 17}\textsuperscript{ème} siècle d'après le mot grec \textit{elektron}, ambre.\\
\textbf{Comment modéliser un phénomène électrostatique ?}

\begin{figure}[h!]
    \begin{minipage}{.49\linewidth}
        \caption{Électrisation par influence, électrisation par contact}
        Selon leur composition, les corps ont des réponses différentes à l'influence de charges électriques :
        \begin{enumerate}
            \item On approche un matériau chargé d'un matériau conducteur : la répartition des charges est modifiée dans le conducteur
            \item On l'approche d'une substance dipolaire ; les dipôles s'orientent
            \item On met en contact un objet chargé avec un objet conducteur : la charge électrique se répartir à la surface
            \item On le met en contact avec un objet non conducteur : une partie de la charge est transférée à l'objet à l'endroit du contact puis les deux objets se repoussent
        \end{enumerate}
    \end{minipage}
    \begin{minipage}{.49\linewidth}
        \caption{Force électrostatique, loi de Coulomb}
        La force électrostatique modélise l'action d'un corps chargé sur un autre. Cette force est décrite par la loi énoncée par Coulomb en 1785 : $$ F_e = k \cdot \dfrac{q_1 \cdot q_2}{d^2}$$
        où $d$ est la distance entre les centres des objets en mètre (\unit{\meter}), $q_1$ et $q_2$ sont les charges électriques portées par les objets 1 et 2 en coulomb (\unit{\coulomb}) et $k = \SI{8,99E9}{\newton\meter\squared\per\coulomb}$.
    \end{minipage}
\end{figure}

\begin{enumerate}
    \item Frotter l'une après l'autre deux pailles contre un morceau de papier essuie-tout. Approcher ensuite les deux pailles l'une de l'autre. Qu'observe-t-on ? Interpréter.
    \item Frotter une baguette de verre contre un morceau de papier essuie-tout. Approcher ensuite la baguette de la paille. Qu'observe-t-on ? Interpréter.
    \item Approcher la paille électrisée d'une petite boule d'aluminium suspendue. qu'observe-t-on dans unn premier temps ? Interpréter. Mettre la paille en contact de la boule d'aluminium. Qu'observe-t-on ?
    \item À l'aide du matériel à disposition, comment dévier sans contact un filet d'eau ?
    \item Expliquer pourquoi dans l'espérience précédente, le filet d'eau subit une force d'attraction et non de répulsion.
    \item Après avoir rappelé l'expression de la force d'interaction gravitatiuonnelle, identifier ses similitdes et ses différences avec la loi de Coulomb.
    \item Expliquer ce qu'il se passe lorsque nos cheveux montent en l'air après avoir enlevé un pull. Comment modélise-t-on cette intéraction ? Plateau, lamelle, tige vesticale.
    \item L'électroscope est constitué d'un plateau métallique relié à une tige verticale métallique fixe. une lamelle métallique peut osciller librement autour d'un axe horizontal lié à cette tige.\\
    Prévoir ce qui va se passer si on électrise par contact le plateau à l'aide d'un bâton d'ébonite frotté et chargé négativement. Consulter l'animation "Electroscope" via le QR-code ci-contre.\\
    Le vérifier expérimentalement. Que se passe-t-il lorsqu'on touche le plateau avec la main ? Quelle propriété du corps humain est ainsi mise en évidence ?
\end{enumerate}
\end{document}